\section{引言}
整除偏序中的密度条件如何保证一条无限链，是 JSP-001022（Erdős Problem \#1217）的核心问题。本文先讨论有限 Abel 求和所提供的密度转移，再说明已有整除链定理的两个组成部分，最后统一处理形式化中的表示和依赖关系。整除链定理的数学来源为 Alexeev 等人的论文\cite{alexeev2026}；本文是证明阐释，不主张新的数论结论。

\subsection{问题陈述与记号}
考虑严格递增的正整数序列 $a_1<a_2<\cdots$，令 $A=\{a_i:i\ge1\}$。对于 $x>e$，定义
\begin{align}
H_A(x)&=\sum_{a\in A,\,a<x}\frac1a,\\
W_A(x)&=\sum_{a\in A,\,2\le a<x}\frac1{a\log a}.
\end{align}
采用自然对数及严格截断 $a<x$。在 $a=1$ 处不计算 $1/(a\log a)$，等价于将该权重定义为零。记
\begin{equation}\label{eq:densities}
\underline\delta_{\log}(A)=\limin\frac{H_A(x)}{\log x},\qquad
\overline\delta_{\log\log}(A)=\limsu\frac{W_A(x)}{\log\log x}.
\end{equation}
以下统一记 $\Delta(A):=\overline\delta_{\log\log}(A)$。调和和的积分比较及第~\ref{sec:density}~节的全体整数加权和估计给出
$0\le\underline\delta_{\log}(A)\le1$ 和 $0\le\Delta(A)\le1$，故这两个密度均为有限实数；后文链计数的归一化上极限则允许为 $+\infty$。对于一个严格递增的指标序列 $k_1<k_2<\cdots$，定义 $b_i=a_{k_i}$ 和
\begin{equation}
C_{a,k}(x)=\#\{i\ge1:a_{k_i}<x\}.
\end{equation}
\begin{theorem}[JSP-001022 的枚举形式]\label{thm:main}
若
\begin{equation}\label{eq:assumption}
\underline\delta_{\log}(A)>0,
\end{equation}
则存在严格递增的正整数指标序列 $(k_i)_{i\ge1}$，满足
\begin{equation}\label{eq:divides}
a_{k_i}\mid a_{k_{i+1}}\qquad(i\ge1),
\end{equation}
并且
\begin{equation}\label{eq:target}
\limsu\frac{C_{a,k}(x)}{\log\log x}
\ge\limsu\frac{W_A(x)}{\log\log x}.
\end{equation}
式 \eqref{eq:target} 按照允许 $+\infty$ 的扩展非负实数上极限解释。
\end{theorem}
应注意量词顺序：只选择\emph{一条固定}的链，而不是令其随截断参数 $x$ 改变。上极限的不等式也不是在每个足够大的 $x$ 逐点成立的结论。


\paragraph{历史文献中的密度表述。}
原文\cite[p.~431]{erdos1966}的文字使用“正下对数密度”，但式~(1) 印为 $\limsup$；第~432~页式~(5) 的结论则比较两个上极限。本文遵循 JSP-001022 目录\cite{jspcatalog}及 Erd\H{o}s Problem \#1217\cite{erdos1217}的下密度表述，以式~\eqref{eq:assumption}~的 $\liminf$ 为假设。这里仅记录原文的文字与公式差异，不将正上密度与正下密度视为等价，也不更改本文主定理。

\subsection{证明路线与本文范围}
证明的主干为
\[
 \underline\delta_{\log}(A)>0
 \ \Longrightarrow\ \overline\delta_{\log\log}(A)>0
 \ \Longrightarrow\ \text{具有定量访问下界的环境整除链}
 \ \Longrightarrow\ \text{集合内无限子链}.
\]
第~\ref{sec:density}~节展开第一步的有限不等式、误差估计和极限过程。第~\ref{sec:chain}~节从不变权重和随机链构造出发，证明访问恒等式、矩估计、确定性路径选择及子链提取。第~\ref{sec:conversion}~节完成端点与枚举转换。最后一节集中给出固定版本源码的对应关系与核查边界。
